\documentclass[DIN, pagenumber=false, fontsize=11pt, parskip=half]{scrartcl}

\usepackage[utf8]{inputenc}
\usepackage[T1]{fontenc}
\usepackage{template}
\usepackage{amsmath}

\DeclareMathOperator{\lcm}{lcm}

\titlehead{\centering\small
    Worked solutions to selected exercises from\\
    Richard Hammack, \emph{Book of Proof, edition 2.2}\\
\rule{0.4\linewidth}{0.4pt}}

\title{}
\author{} % Oscar A. Carrera
\date{}

\begin{document}
\maketitle

\setcounter{section}{9}

\section{Mathematical Induction}

\textbf{Prove the following statements with either induction, strong induction or proof by smallest counterexample}.

\begin{problems}
    \problem For every integer $n \in \mathbb{N}$, it follows that
    $1^2 + 2^2 + 3^2 + 4^2 + \mathellipsis + n^2 = \dfrac{n(n+1)(2n+1)}{6}$.
    \begin{proof}
        We prove the statement by mathematical induction.

        \begin{enumerate}
            \item Observe that if $n = 1$, this statement is $ 1^2 = \dfrac{(1)(2)(3)}{6} = \dfrac{6}{6}$, which is true.

            \item Assume $1^2 + 2^2 + 3^2 + 4^2 + \mathellipsis k^2 = \dfrac{k(k+1)(2k+1)}{6}$ for some integer $n = k \geq 1$. Observe that this implies the statement is true for $n = k + 1$.
                \begin{align*}
                    1^2 + 2^2 + 3^2 + \mathellipsis + k^2 + (k+1)^2 & = \\ 
                    (1^2 + 2^2 +3^2 + \mathellipsis + k^2) + (k+1)^2 & = \\ 
                    \dfrac{k(k+1)(2k+1)}{6} + (k+1)^2 & = \dfrac{k(k+1)(2k+1)}{6} + \dfrac{6(k+1)^2}{6} \\ 
                                                      & = \dfrac{k(k+1)(2k+1)+6(k+1)(k+1)}{6} \\ 
                                                      & = \dfrac{2k^3+9k^2+13k+6}{6}\\ 
                                                      & = \dfrac{(k+1)(k+2)(2k+3)}{6}\\ 
                                                      & = \dfrac{(k+1)([k+1] + 1)(2[k+1]+1)}{6}
                \end{align*}
        \end{enumerate}
        Therefore, $1^2 + 2^2 +3^3 + \mathellipsis + k^2 + (k+1)^2 = \dfrac{(k+1)([k+1] + 1)(2[k+1]+1)}{6}$, which means the statement is true for $n = k + 1$.
    \end{proof}

    \problem If $n \in \mathbb{N}$, then $1\cdot2 + 2\cdot3 + 3\cdot4 + 4\cdot5 + \mathellipsis + n(n+1) = \dfrac{n(n+1)(n+2)}{3}$.
    \begin{proof}
        We prove the statement by mathematical induction.

        \begin{enumerate}
            \item Observe that if $n=1$, this statement is $1(1+1) = \dfrac{1(1+1)(1+2)}{3} = 2$, which is true.
                
            \item Assume $1\cdot2 + 2\cdot3 + 3\cdot4 + \mathellipsis + k(k+1) = \dfrac{k(k+1)(k+2)}{3}$ for some integer $n = k \geq 1$. Observe that this implies the statement is true for $n = k + 1$.
                \begin{align*}
                    1\cdot2 + 2\cdot3 + \mathellipsis + k(k+1) + (k+1)(k+2) & = \\ 
                    (1\cdot2 + 2\cdot3  \mathellipsis + k(k+1)) + (k+1)(k+2) & = \\ 
                    \dfrac{k(k+1)(k+2)}{3} + (k+1)(k+2) & = \dfrac{k(k+1)(k+2)}{3} + \dfrac{3(k+1)(k+2)}{3}\\ 
                                                        & = \dfrac{k(k^2+3k+2) + 3(k^2+3k+2)}{3}\\ 
                                                        & = \dfrac{k^3+6k^2+11k+6}{3}\\ 
                                                        & = \dfrac{(k+1)(k^2+5k+6)}{3}\\ 
                                                        & = \dfrac{(k+1)(k+2)(k+3)}{3}\\ 
                                                        & = \dfrac{(k+1)([k+1]+1)([k+1]+2)}{3}
                \end{align*}
        \end{enumerate}
        Therefore $1\cdot2+2\cdot3+3\cdot4+\mathellipsis+k(k+1) + (k+1)(k+2) = \dfrac{(k+1)([k+1]+1)([k+1]+2)}{3}$, which means the statement is true for $n = k + 1$.
    \end{proof}
\end{problems}

\end{document}
